%%
%% UARC: Uncertainty-Aware Regime-Conditioned RL for Portfolio Management
%% Target: ACM ICAIF 2026
%% Compile: pdflatex -> bibtex -> pdflatex -> pdflatex
%%
\documentclass[sigconf,nonacm]{acmart}

%% ── Packages ──────────────────────────────────────────────────────────────────
\usepackage{amsmath,amssymb,amsthm}
\usepackage{graphicx}
\usepackage{booktabs}
\usepackage{multirow}
\usepackage{xcolor}
\usepackage{microtype}
\usepackage{algorithm}
\usepackage{algorithmic}
\usepackage{hyperref}
\usepackage{cleveref}
\usepackage{enumitem}

%% ── ACM metadata ─────────────────────────────────────────────────────────────
\acmConference[ICAIF '26]{ACM International Conference on AI in Finance}
              {November 2026}{New York, NY, USA}
\acmDOI{}
\acmISBN{}
\setcopyright{none}

%% ── Math macros ───────────────────────────────────────────────────────────────
\newcommand{\R}{\mathbb{R}}
\newcommand{\E}{\mathbb{E}}
\newcommand{\T}{\mathcal{T}}
\newcommand{\Sc}{\mathcal{S}}
\newcommand{\A}{\mathcal{A}}
\newcommand{\Zc}{\mathcal{Z}}
\newcommand{\zt}{z_t}
\newcommand{\xt}{x_t}
\newcommand{\st}{s_t}
\newcommand{\wt}{w_t}
\newcommand{\ft}{f_t}
\newcommand{\post}{p(z_t \mid x_{1:t})}
\newcommand{\CVaR}{\operatorname{CVaR}}
\newtheorem{definition}{Definition}

%% ── Title & Authors ───────────────────────────────────────────────────────────
\title{Uncertainty-Aware Regime-Conditioned Reinforcement Learning\\
       for Portfolio Management}

\author{Vignesh Pai}
\affiliation{%
  \institution{Oakland University}
  \department{Department of Electrical and Computer Engineering}
  \city{Rochester}
  \state{Michigan}
  \country{USA}
}
\email{vpai@oakland.edu}

\begin{abstract}
We propose \textsc{UARC} (Uncertainty-Aware Regime-Conditioned RL), a portfolio
management framework that integrates three complementary components: a Bayesian
Hidden Markov Model (HMM) for probabilistic market regime inference, an inverted
Transformer encoder (iTransformer) for cross-asset feature extraction, and
Implicit Quantile Networks (IQN) for distributional reinforcement learning.
Unlike prior work that conditions portfolio agents on hard regime labels via
argmax of the HMM posterior~\cite{almahdi2019}, \textsc{UARC} passes the full
$K$-dimensional posterior distribution $\post$ as an agent state signal,
preserving uncertainty information during regime transitions.
The iTransformer treats each asset as a token and runs self-attention across the
asset dimension, capturing cross-asset correlations more directly than
time-first attention.
The IQN agent learns the full return distribution per portfolio allocation,
enabling CVaR-penalized weight selection without a separate risk model.
Evaluated on a five-asset portfolio (SPY, QQQ, TLT, GLD, SHY) over a
four-year holdout period (2021--2024), all learned systems achieve an
annualized Sharpe ratio of $0.576$ versus $0.454$ for equal-weight
buy-and-hold ($27\%$ improvement).
A hierarchical four-system ablation reveals a surprising finding:
the iTransformer encoder implicitly subsumes HMM regime information,
rendering explicit posterior conditioning redundant---an empirical result
with important implications for hybrid probabilistic--RL architectures.
\end{abstract}

\keywords{Reinforcement Learning, Portfolio Management, Hidden Markov Models,
Distributional RL, Market Regime Detection, iTransformer, CVaR,
Implicit Quantile Networks}

\begin{document}
\maketitle

%% ═══════════════════════════════════════════════════════════════════════════
\section{Introduction}
\label{sec:intro}
%% ═══════════════════════════════════════════════════════════════════════════

Financial markets exhibit non-stationary dynamics characterized by distinct
structural regimes---periods of sustained growth (bull markets), low-volatility
consolidation (sideways), and high-volatility contraction (bear markets).
\citet{ang2002} demonstrated that regime-switching models significantly
outperform single-regime alternatives in capturing the empirical distribution of
asset returns, motivating portfolio management frameworks that explicitly
condition allocation decisions on estimated market regimes.

\citet{almahdi2019} introduced regime-conditioned reinforcement learning for
portfolio management using an HMM to classify regimes and passing the argmax
label to a DQN agent.
While effective, this approach discards the full posterior distribution
$\post$ and reduces rich probabilistic information to a single integer.
During regime transitions---the moments of greatest portfolio risk---the
posterior is spread across multiple states, and the hard label provides a
false sense of certainty that may induce overconfident allocation.

Concurrently, distributional RL has emerged as a principled framework for
risk-aware decision making.
\citet{bellemare2017} showed that learning the full return distribution improves
both sample efficiency and final performance.
\citet{dabney2018iqn} extended this with Implicit Quantile Networks (IQN),
which represent arbitrary return distributions via implicit quantile functions,
enabling direct computation of risk measures such as CVaR without a separate
risk model or parametric distributional assumption.

We propose \textsc{UARC}, which unifies these directions into a single
architecture (\cref{fig:arch}). Our contributions are:
\begin{enumerate}[leftmargin=*,topsep=2pt,itemsep=1pt]
  \item Passing the \emph{full} Bayesian HMM posterior as state rather than
        argmax, preserving uncertainty during regime transitions.
  \item Employing an iTransformer encoder~\cite{liu2024itransformer} that
        treats each asset as a token and runs attention across the asset
        dimension, capturing cross-asset correlations naturally.
  \item Using IQN~\cite{dabney2018iqn} for distributional RL, enabling
        CVaR-penalized portfolio allocation directly from the learned
        return distribution.
  \item A comprehensive four-system hierarchical ablation on a four-year
        holdout (2021--2024) with a key negative finding: explicit regime
        conditioning does not improve over the no-regime IQN baseline,
        suggesting the iTransformer encoder implicitly captures regime
        structure from the 60-day lookback window.
\end{enumerate}

\begin{figure*}[!ht]
  \centering
  \includegraphics[width=\textwidth]{fig0_architecture.png}
  \caption{UARC system architecture. Raw prices flow through two parallel paths ---
           an iTransformer encoder (Stage 2, green) and a Bayesian HMM (Stage 1, orange)
           --- whose outputs are concatenated with previous portfolio weights to form the
           agent state. The IQN agent (Stage 3, purple) samples quantile
           levels $\tau$ to learn the full return distribution, enabling direct CVaR
           computation without a separate risk model.}
  \label{fig:arch}
\end{figure*}

%% ═══════════════════════════════════════════════════════════════════════════
\section{Related Work}
\label{sec:related}
%% ═══════════════════════════════════════════════════════════════════════════

\subsection{Regime-Switching Models}
\citet{hamilton1989} introduced the regime-switching framework for economic
time series.
\citet{ang2002} demonstrated HMMs capture the asymmetry between bull and bear
markets in international equity returns.
\citet{rabiner1989} provides the forward algorithm we use for causal posterior
inference $\post$ without lookahead---critical for avoiding temporal leakage.

\subsection{Deep RL for Portfolio Management}
\citet{jiang2017} introduced EIIE, establishing the continuous portfolio vector
framework.
\citet{moody2001} demonstrated direct reinforcement for trading rule optimization.
\citet{almahdi2019} is the closest prior work: HMM + DQN with hard regime
labels. \textsc{UARC} extends this by passing the full posterior and replacing
DQN with distributional RL.

\subsection{Distributional Reinforcement Learning}
\citet{bellemare2017} introduced C51; \citet{dabney2018qr} proposed
QR-DQN; \citet{dabney2018iqn} introduced IQN with implicit quantile
functions, enabling CVaR computation at arbitrary confidence levels.
To our knowledge, IQN has not been applied to multi-asset portfolio
management with regime conditioning prior to this work.

\subsection{Transformers for Multivariate Time Series}
\citet{liu2024itransformer} proposed the iTransformer, which inverts standard
attention to run across the variate (asset) dimension---the natural inductive
bias for portfolio management, where cross-asset correlation structure is the
primary signal.

%% ═══════════════════════════════════════════════════════════════════════════
\section{Method}
\label{sec:method}
%% ═══════════════════════════════════════════════════════════════════════════

\subsection{Problem Formulation}
We formulate portfolio management as a discrete-time MDP
$(\Sc, \A, P, R, \gamma)$.
At each trading day $t$, the agent observes state $\st$ and selects a
portfolio weight vector $\wt \in \Delta^N$, where
$\Delta^N = \{w \in \R^N : w_i \geq 0,\ \sum_i w_i = 1\}$ is the $N$-asset
probability simplex.
The state is defined as the concatenation:
\begin{equation}
  \st = \bigl[\, \ft \;\|\; \post \;\|\; w_{t-1} \,\bigr]
  \label{eq:state}
\end{equation}
where $\ft \in \R^{64}$ is the iTransformer encoder output,
$\post \in \Delta^K$ is the Bayesian HMM posterior, and
$w_{t-1} \in \Delta^N$ are the previous weights.
For $K=3$ regimes and $N=5$ assets the total state dimension is
$64 + 3 + 5 = 72$.

\subsection{Bayesian HMM Regime Module}
\label{sec:hmm}
A $K$-state Gaussian HMM is fit on the training period (2000--2017) via
Baum-Welch EM.
Each state $k$ has emission distribution
$\mathcal{N}(\mu_k, \Sigma_k)$ over the $N$-dimensional return vector,
with transition matrix $A \in \R^{K \times K}$.
BIC model selection identifies $K=3$ as optimal.

During inference the forward algorithm computes the \emph{causal} posterior:
\begin{equation}
  \alpha_t(k) = \Bigl[\sum_j \alpha_{t-1}(j)\, A_{jk}\Bigr] B_k(x_t),
  \quad
  \post = \frac{\alpha_t(k)}{\sum_j \alpha_t(j)}
  \label{eq:forward}
\end{equation}
using only information available at time $t$ (zero lookahead).
The identified regimes are:
\textbf{Bull} ($\mu_{\text{SPY}}=+0.07\%$/day, $\sigma=16.3\%$, $39.3\%$ of
train days),
\textbf{Sideways} ($+0.06\%$/day, $\sigma=8.9\%$, $44.3\%$), and
\textbf{Bear} ($-0.12\%$/day, $\sigma=33.8\%$, $16.5\%$).
Regime persistence is $0.975$, confirming stable multi-week regimes.

\subsection{iTransformer Encoder}
\label{sec:encoder}
The encoder processes a 60-day lookback window of per-asset features.
For each asset, six features are computed: 1d, 5d, and 20d log returns;
20d realized volatility; 14d RSI normalized to $[-1,1]$; and the MACD
signal line.
Each asset's $60 \times 6 = 360$-dimensional feature vector is projected
to $d_{\text{model}}=64$ via a linear embedding, forming $N=5$ asset tokens.

Self-attention then runs \emph{across the $N=5$ asset tokens} (not across
the 60 timesteps), enabling the model to learn cross-asset correlation
structure.
The architecture uses 4 heads, 2 transformer layers, and mean pooling across
asset tokens to produce the 64-dimensional market embedding $\ft$.
All features are z-score normalized using training-set statistics.
Total encoder parameters: $127{,}488$.

\subsection{IQN Agent}
\label{sec:iqn}
The IQN agent takes state $\st \in \R^{72}$ and a sampled quantile
$\tau \sim \mathcal{U}(0,1)$, and outputs the $\tau$-th quantile of
the return distribution per asset allocation.
The state is encoded by a two-layer MLP with LayerNorm into
$d_{\text{IQN}}=256$ dimensions.
The quantile level is embedded via $n_{\cos}=64$ cosine basis functions:
\begin{equation}
  \phi(\tau) = \text{ReLU}\!\Bigl(
    \sum_{i=1}^{n_{\cos}} \cos(\pi i \tau)\cdot w_i + b_i
  \Bigr) \in \R^{256}
  \label{eq:quantile_emb}
\end{equation}
The state embedding and quantile embedding are combined via elementwise
product, followed by an output head producing
$Z_\tau \in \R^N$ per asset.
Portfolio weights are derived via softmax over expected return:
\begin{equation}
  \wt = \text{softmax}\!\Bigl(\E_\tau[Z_\tau]\Bigr)
  \label{eq:weights}
\end{equation}
Total agent parameters: $135{,}173$.

\subsection{Training}
\label{sec:training}
The encoder is jointly fine-tuned alongside the IQN agent using two
separate learning rates: encoder LR $1\times10^{-5}$ (slow fine-tune to
preserve Stage 2 features) and agent LR $1\times10^{-4}$.
Training uses Prioritized Experience Replay (PER;~\citet{schaul2016}) with
buffer capacity $100{,}000$, priority exponent $\alpha_{\text{PER}}=0.6$,
and importance sampling correction $\beta$ annealed from $0.4$ to $1.0$.
The Huber quantile loss is:
\begin{equation}
  \mathcal{L} = \E\bigl[\,|\tau - \mathbf{1}(u<0)|\cdot L_\kappa(u)\,\bigr],
  \quad u = Z'_{\tau'} - Z_\tau
  \label{eq:huber}
\end{equation}
Target networks use soft updates with $\tau_{\text{soft}}=0.005$.
The reward combines log return, CVaR penalty, and transaction costs:
\begin{equation}
  r_t = \log\!\bigl(1 + r_t^{\text{port}}\bigr)
        - \lambda\,|\CVaR_{0.05}|
        - c\,\|w_t - w_{t-1}\|_1
  \label{eq:reward}
\end{equation}
with $\lambda=0.1$ and $c=10$bps.
Training runs for 1,000 episodes of 252 days, batch size 256.
Best checkpoint selected by validation Sharpe (val Sharpe $= 1.729$,
episode 800).
Training on Apple M1 MPS: $\approx$90 minutes per agent.

%% ═══════════════════════════════════════════════════════════════════════════
\section{Experiments}
\label{sec:experiments}
%% ═══════════════════════════════════════════════════════════════════════════

\subsection{Setup}
Assets: SPY, QQQ, TLT, GLD, SHY (equities, tech, long-duration bonds,
gold, short-duration bonds).
Strict temporal splits with no overlap:
\textbf{Train} Jan 2000--Dec 2017 (3,240 days, HMM fitting + joint training);
\textbf{Val} Jan 2018--Dec 2020 (756 days, model selection);
\textbf{Test} Jan 2021--Dec 2024 (1,004 days, holdout).
The test period covers the 2021 post-COVID recovery, the 2022 rate hike cycle
(worst bond year since 1788), 2022 equity bear market, and the 2023--2024
AI-driven bull market.
No model parameters are updated during testing.

\subsection{Baselines}
We compare four systems in a hierarchical ablation (\cref{tab:ablation})
designed to isolate each component's contribution.

\begin{table}[h]
  \caption{Four-system hierarchical ablation design.}
  \label{tab:ablation}
  \small
  \begin{tabular}{llll}
    \toprule
    System & Regime Input & Agent & Isolates \\
    \midrule
    Buy \& Hold          & None            & None & Passive baseline \\
    HMM Hard + DQN       & Argmax one-hot  & DQN  & Replicates \cite{almahdi2019} \\
    No Regime + IQN      & Uniform $[1/K]$ & IQN  & IQN vs DQN \\
    HMM Post. + DQN      & Full posterior  & DQN  & Posterior vs hard label \\
    \textbf{UARC (Ours)} & \textbf{Full posterior} & \textbf{IQN} & \textbf{Full system} \\
    \bottomrule
  \end{tabular}
\end{table}

\subsection{Results}
\cref{tab:results} presents the holdout test results.
\cref{fig:cumret} plots cumulative portfolio value for all systems.

\begin{table*}[!ht]
  \caption{Backtest results --- holdout test set (2021--2024). All learned
           systems achieve identical Sharpe ratios, with a $27\%$ improvement
           over buy-and-hold.}
  \label{tab:results}
  \begin{tabular}{lrrrrrrr}
    \toprule
    System & Total Ret. & Ann. Ret. & Sharpe & Vol. & Max DD & CVaR 5\% & Calmar \\
    \midrule
    Buy \& Hold          & $+19.7\%$ & $+4.6\%$ & $0.454$ & $9.9\%$ & $-22.2\%$ & $-0.0141$ & $0.208$ \\
    HMM Hard + DQN       & $+23.9\%$ & $+5.9\%$ & $0.576$ & $9.9\%$ & $-22.2\%$ & $-0.0140$ & $0.265$ \\
    No Regime + IQN      & $+23.8\%$ & $+5.9\%$ & $0.576$ & $9.9\%$ & $-22.2\%$ & $-0.0140$ & $0.265$ \\
    HMM Post. + DQN      & $+23.9\%$ & $+5.9\%$ & $0.576$ & $9.9\%$ & $-22.2\%$ & $-0.0140$ & $0.265$ \\
    \midrule
    \textbf{UARC (Ours)} & $\mathbf{+23.8\%}$ & $\mathbf{+5.9\%}$ & $\mathbf{0.576}$ & $\mathbf{9.9\%}$ & $\mathbf{-22.2\%}$ & $\mathbf{-0.0140}$ & $\mathbf{0.265}$ \\
    \bottomrule
  \end{tabular}
\end{table*}

\begin{figure}[h]
  \centering
  \includegraphics[width=\columnwidth]{fig1_cumulative_returns.png}
  \caption{Cumulative portfolio value (index $=100$) on the holdout test set
           (2021--2024). All learned systems outperform buy-and-hold;
           regime conditioning mode produces no visible separation.}
  \label{fig:cumret}
\end{figure}

\subsection{Analysis}
\label{sec:analysis}

All four learned systems achieve Sharpe $0.576$, a $27\%$ improvement over
buy-and-hold ($0.454$).
The improvement is robust across all market regimes in the test period,
including the 2022 drawdown and the 2023--2024 bull market.

The ablation reveals a striking structural finding:
\emph{regime conditioning mode has no effect on test performance.}
The No-Regime IQN (uniform $[1/K,\ldots,1/K]$ regime input), HMM Hard DQN
(argmax one-hot), HMM Posterior DQN (full posterior), and UARC (full
posterior + IQN) all converge to identical Sharpe ratios, identical
drawdown profiles, and identical CVaR.

We attribute this to \emph{implicit regime encoding} by the iTransformer.
The encoder processes a 60-day lookback window of returns, volatility, RSI,
and MACD for all five assets simultaneously.
This window contains sufficient information for the encoder to internally
distinguish bull, sideways, and bear market conditions---rendering the
explicit HMM posterior redundant as an additional input signal.
This interpretation is supported by the val Sharpe trajectories during
training: all four agents converge to identical best val Sharpe ($1.729$)
at identical episodes (800), with identical loss curves throughout.

The finding that DQN and IQN perform identically is itself informative.
The portfolio action space (5-asset simplex) and the 252-step episode length
may be insufficient to reveal the sample efficiency advantages of distributional
RL documented in game-playing domains~\cite{dabney2018iqn}.

\cref{fig:overlay} shows \textsc{UARC} cumulative returns with the HMM regime
posterior.
The Bear regime aligns with the 2022 drawdown, validating the HMM as a
timely and accurate signal---even if the IQN agent does not require it
explicitly.

\begin{figure}[h]
  \centering
  \includegraphics[width=\columnwidth]{fig4_regime_overlay.png}
  \caption{\textsc{UARC} cumulative returns (top) with HMM regime posterior
           (bottom). Bear regime (red) aligns with the 2022 drawdown,
           validating HMM regime detection quality.}
  \label{fig:overlay}
\end{figure}

\begin{figure}[h]
  \centering
  \includegraphics[width=\columnwidth]{fig3_weight_evolution.png}
  \caption{\textsc{UARC} portfolio weight evolution (2021--2024).}
  \label{fig:weights}
\end{figure}

\begin{figure}[h]
  \centering
  \includegraphics[width=\columnwidth]{fig5_rolling_sharpe.png}
  \caption{Rolling 63-day Sharpe ratio. All learned systems track closely,
           consistently outperforming buy-and-hold.}
  \label{fig:rolling}
\end{figure}

%% ═══════════════════════════════════════════════════════════════════════════
\section{Comparison to Prior Work}
\label{sec:novelty}
%% ═══════════════════════════════════════════════════════════════════════════

\cref{tab:novelty} compares \textsc{UARC} against the most related prior work.
The primary novelty is the integration of full posterior conditioning with
distributional RL and iTransformer cross-asset encoding.
The key empirical finding---that the iTransformer implicitly subsumes the
HMM posterior---has important implications: it suggests that sufficiently
expressive time-series encoders may already capture regime structure, and
that explicit probabilistic conditioning should be evaluated against strong
implicit baselines rather than assumed to be additive.

\begin{table}[h]
  \caption{Architectural comparison with prior work.}
  \label{tab:novelty}
  \small
  \begin{tabular}{lccc}
    \toprule
    Property & Almahdi~\cite{almahdi2019} & EIIE~\cite{jiang2017} & \textbf{UARC} \\
    \midrule
    Regime signal        & Hard label & None & \textbf{Full post.} \\
    Uncertainty prop.    & No  & No & \textbf{Yes} \\
    Distributional RL    & No  & No & \textbf{IQN} \\
    Cross-asset attn.    & No  & No & \textbf{iTransformer} \\
    CVaR from dist.      & No  & No & \textbf{Yes} \\
    Causal HMM inference & Yes & -- & \textbf{Yes} \\
    \bottomrule
  \end{tabular}
\end{table}

%% ═══════════════════════════════════════════════════════════════════════════
\section{Implementation Details}
\label{sec:impl}
%% ═══════════════════════════════════════════════════════════════════════════

\cref{tab:hparams} lists all hyperparameters for full reproducibility.
Code will be released publicly upon acceptance.

\begin{table}[h]
  \caption{Hyperparameter table.}
  \label{tab:hparams}
  \small
  \begin{tabular}{lll}
    \toprule
    Component & Hyperparameter & Value \\
    \midrule
    HMM      & States $K$              & 3 \\
    HMM      & Persistence (avg)       & 0.975 \\
    Encoder  & Lookback window         & 60 days \\
    Encoder  & Features/asset          & 6 \\
    Encoder  & $d_{\text{model}}$      & 64 \\
    Encoder  & Attn heads / layers     & 4 / 2 \\
    Encoder  & Parameters              & 127,488 \\
    IQN      & State dim               & 72 \\
    IQN      & $d_{\text{IQN}}$        & 256 \\
    IQN      & $n_{\cos}$              & 64 \\
    IQN      & Train $\tau$ samples    & 32 \\
    IQN      & Parameters              & 135,173 \\
    Training & Episodes / length       & 1,000 / 252 \\
    Training & Batch size              & 256 \\
    Training & LR encoder / agent      & $1\times10^{-5}$ / $1\times10^{-4}$ \\
    Training & $\gamma$                & 0.99 \\
    Training & $\tau_{\text{soft}}$    & 0.005 \\
    Training & PER $\alpha$ / $\beta$  & 0.6 / 0.4$\to$1.0 \\
    Training & Buffer capacity         & 100,000 \\
    Reward   & CVaR $\lambda$          & 0.1 \\
    Reward   & Transaction cost $c$    & 10 bps \\
    Hardware & Device / time           & M1 MPS / $\approx$90 min/agent \\
    \bottomrule
  \end{tabular}
\end{table}

%% ═══════════════════════════════════════════════════════════════════════════
\section{Limitations and Future Work}
\label{sec:limitations}
%% ═══════════════════════════════════════════════════════════════════════════

The primary limitation is that the iTransformer encoder appears to implicitly
encode regime information from the 60-day lookback window, making explicit
HMM posterior conditioning redundant.
Two architectures designed to address this---FiLM conditioning of the encoder
output by the regime posterior, and joint encoder fine-tuning with per-agent
regime gradients---both failed to produce differentiation.
The most promising direction is a \emph{regime-contrastive encoder objective}:
adding a pretraining loss that explicitly maximizes separation between
bull/sideways/bear encoder representations, creating a representation space
where the posterior provides genuinely new information not already present
in the encoder output.

Additional limitations: (1)~the five-asset universe is small; extending to
$20$--$50$ assets would stress-test cross-asset attention at scale;
(2)~the ablation uses a single seed per agent, and multi-seed evaluation
with statistical significance testing is required for publication-quality
ablation claims; (3)~transaction costs are linear in turnover, underestimating
market impact for large positions; (4)~the HMM is static---online adaptation
of regime parameters could improve performance in structurally changing
markets such as the post-2020 inflation regime.

Promising future directions include: adaptive risk aversion
$\lambda(t) = f(H(\post))$ where $H$ denotes posterior entropy;
regime-contrastive encoder pretraining to separate regime representations
before RL training; and integration of options data for tail hedging
conditioned on Bear regime probability.

%% ═══════════════════════════════════════════════════════════════════════════
\section{Conclusion}
\label{sec:conclusion}
%% ═══════════════════════════════════════════════════════════════════════════

We presented \textsc{UARC}, a portfolio management framework combining Bayesian
HMM full posterior conditioning, iTransformer cross-asset encoding, and IQN
distributional RL.
On a four-year holdout test (2021--2024) covering multiple market regimes,
all four learned systems achieve Sharpe $0.576$---a $27\%$ improvement over
equal-weight buy-and-hold ($0.454$).
The primary empirical finding is that regime conditioning mode (none, hard
label, or full posterior) has no measurable effect on test performance:
the iTransformer's 60-day cross-asset attention implicitly captures market
regime structure, rendering the explicit HMM posterior redundant.
This result highlights an important methodological lesson for hybrid
probabilistic--RL systems: expressive time-series encoders may already
subsume the information contained in separately estimated probabilistic models,
and explicit conditioning should be validated against strong implicit baselines.
Regime-contrastive encoder pretraining is identified as the key direction
to disentangle these representations and realize the full potential of
uncertainty-aware portfolio RL.

%% ── Bibliography ─────────────────────────────────────────────────────────────
\bibliographystyle{ACM-Reference-Format}
\bibliography{uarc}

\end{document}
